\chapter{Researching a person}
\label{ch:person}

\section*{What this chapter is for}

At this level the people interviewing you have usually written something in
public. Reading it takes an hour and changes the conversation more reliably
than another hour of technical revision, because it lets you calibrate to a
person rather than to a role.

\section{Why it works}

An interview is a conversation with somebody who has opinions about the subject
and has usually published them. Knowing what they think does three things:

\textbf{It tells you what they will probe.} Somebody who has written about
evaluation will ask about evaluation, and will not be satisfied by the version
of the answer that satisfies a generalist.

\textbf{It tells you where agreement is cheap and where it is not.} If you
agree with a published position, saying so with a reason is a real exchange. If
you disagree, knowing that in advance means you can disagree well rather than
stumbling into it.

\textbf{It gives you questions worth asking.} A question that could only be
asked of that person is worth ten generic ones, and it is the clearest
available signal that you prepared for them specifically.

\judgement{The effect is asymmetric in your favour: they have read your
application once, and you can have read a year of their thinking. That is not a
trick, it is just an available asymmetry, and not using it is the odd choice.}

\section{Where people at this level publish}

In roughly descending order of yield:

\begin{itemize}
  \item A personal site or blog. The highest-signal source and the most often
        missed, because it is not on the professional network.
  \item Long-form posts on engineering publications and aggregators.
  \item Conference talks and their slides, which frequently contain the
        architecture diagram you want.
  \item Open-source contributions: what they maintain, what they review, what
        they argue about in issues.
  \item Academic publications, where the person came through research.
  \item The company's own engineering blog, where their name is on a post.
\end{itemize}

\begin{kitnote}
Read what they wrote, not what a summary says about it. A position summarised
is a position you will misquote, and misquoting somebody's own writing back to
them is worse than not having read it.
\end{kitnote}

\section{What to do with it}

\begin{kitmethod}
\textbf{Find the bridges.} Two or three genuine points of contact between what
they have written and what you have done. Genuine is the constraint: a bridge
you manufactured will not survive a follow-up question.

\textbf{Write one line per bridge}, in your words, including what you would say
if they push back. Not a script \dash{} a position you can hold.

\textbf{Write one question that only they could answer.} About a specific thing
they wrote, and a genuine question rather than a compliment in question form.

\textbf{Write down the disagreement, if there is one.} If you think a published
position is wrong, decide in advance whether to raise it and how. Usually the
answer is to raise it as a question about conditions \dash{} ``does that still
hold when \ldots'' \dash{} rather than as a contradiction.
\end{kitmethod}

\code{templates/interviewer-brief.md} is one page per person and takes twenty
minutes once you have found the material.

\section{The line}

Research the professional record. Do not research the person.

\begin{kittrap}
The boundary is not subtle, and it is worth stating because the tooling makes
crossing it easy. What somebody has published about their work is offered to
you. What can be assembled about where they live, who they know, or what they
were doing five years ago is not, and mentioning it converts a well-prepared
candidate into an unsettling one in a single sentence.

The test: could you say how you know this, out loud, without it being strange?
``I read your post on evaluation'' passes. Anything that would require an
explanation does not.
\end{kittrap}

\section{When there is nothing to find}

Some people publish nothing, and that is common. Fall back to the layer above:
what the company publishes, what the team's open-source output looks like, what
their job postings emphasise. Then go into the conversation prepared to find
out in the first five minutes what they care about \dash{} which is what you
were trying to learn in advance anyway.

\judgement{A candidate who asks, early and genuinely, what the hardest problem
in the team currently is will often get the calibration they were looking for
in the first two minutes of the call.}

\begin{drill}
For the next person you are scheduled to speak to, fill in
\code{templates/interviewer-brief.md}: what they have published, two bridges to
your own work, one question only they could answer, and any disagreement you
would need to handle. Twenty minutes. If you find nothing, write that down too
\dash{} it changes your opening rather than leaving you to discover it live.
\end{drill}

\section*{What this chapter does not claim}

It does not claim researching an interviewer changes the outcome. It changes
the conversation, which is a smaller and more defensible claim, and the
mechanism is straightforward: you spend less of the call working out what they
care about.

The yield ordering of the sources is the author's judgement from doing this,
not a measurement.
